\section{相关工作与研究边界}\label{sec:related}

跨技术通信按作用层次可分为物理层仿真与包级/多载波并发广播两类。在物理层方向，BlueBee~\cite{jiang2017bluebee} 利用 \SI{1}{Mbps} \ble{} 高斯频移键控（GFSK）波形构造 \stdfour{} 合法码片序列，每字节 \stdfour{} 帧需占用 $4$ 字节 \ble{} 载荷，并在商用芯片上实现了双标准兼容帧解析与智能灯具控制；其期刊扩展版进一步面向商用物联网设备提升了跨技术通信比特率~\cite{jiang2019boosting}。在包级与多协议广播方向，Brunner 等人~\cite{brunner2021leveraging} 基于包级能量突发（X-Burst）与原生协议栈共存调度，实现了由智能手机无网关控制 \ble{} 门锁与 ZigBee 灯具的异构智能家居原型；WiCast~\cite{xia2025wicast} 则在 IEEE~802.11ax（Wi-Fi~6）正交频分多址（OFDMA）单个资源单元上合成复合信号，并行服务 \ble{}、ZigBee 与 LoRa 接收端。与上述工作不同，本文聚焦于单个 \SI{2}{MHz} 带宽 \ble{}~5 扩展广播发送机约束，既不依赖宽带 Wi-Fi OFDMA 发射端，也不把双协议广播本身视作首次提出，而是定量评估在单个 \texttt{AUX\_ADV\_IND} 中并置双侧结构化控制记录与时间元数据时的布局开销、波形成本与联合解析权衡。

在跨技术时间分发方向，专用同步协议致力于实现异构节点间亚毫秒乃至微秒级绝对时钟对齐。X-Sync~\cite{hofmann2021xsync} 在商用 \ble{} 与 \stdfour{} 节点间采用变长能量突发与重复接收信号强度（RSS）二分搜索，借助外部参考测得微秒级跨设备时钟偏差；Crocs~\cite{yu2018crocs} 与 TimeBee~\cite{gao2023timebee} 面向 Wi-Fi 与 ZigBee 共存网络设计了跨技术时间同步框架。本文并不提出新的跨设备时钟同步算法，也不将本系统的测量值与 X-Sync 等同步精度指标直接排名：本文将共享时隙参考元数据附属嵌入单向共包控制流，仅在固定次信道下评估单侧接收机本侧 IQ 采样域的留出预测残差、未知外部参考界下的保守拒绝以及单侧失更后的状态恢复行为。

\ble{}~5.1 核心规范定义的扩展广播机制~\cite{bluetooth_core51} 将主信道指针与次信道长载荷解耦，突破了传统广播 \SI{31}{B} 荷载限制。为保证实验归因清晰并遵守测量边界，本文在设计阶段明确以下范围与非目标：确认性实验在固定次信道上开展（不将动态 \texttt{AuxPtr} 切频跟随瞬态与在线时间机制混同），其中 \modejc{} 与 \modesc{} 的 \ble{} 侧均使用 ch3（\SI{2410}{MHz}）、\stdfour{} 侧均使用 ch12（\SI{2410}{MHz}），构成同频可协议归因对照；而 \modelc{} 的 \ble{} 侧位于传统广播信道 ch39（\SI{2480}{MHz}），因此 \modejc{}--\modelc{} 的联合解析率差异仅对应所测封装与信道组合。此外，所有成本均以提交设计波形代理时长（\si{\micro\second}）核算，不宣称功率计实测射频空口时长或能耗；全部结果基于单 SDR 发送机与两套独立 HackRF 接收链，不宣称商用协议栈完整兼容或物理执行器同步动作。
